% =========================================================================
% §3 System Description (SuperCoder + Context Engine)
% Budget: ~2.5 pp
% Drafting order: After Results (so we know which mechanics matter).
%
% Content blueprint:
%   - 3.1 SuperCoder harness; agent loop, tools, modes (Ask/Plan/Coding),
%     turn structure.
%   - 3.2 Context engine (what gets ablated); index construction
%     (tree-sitter parse, call-graph, embeddings, BM25); agent-facing tools
%     (codebase_search, codebase_graph); what each returns.
%     ON vs OFF concretely: ON = both tools available; OFF = tools removed
%     from the toolset, all else identical.
%   - 3.3 Comparator: OpenCode (half-page); agent loop, tools, retrieval
%     approach (ripgrep + read), what's not present (no persistent index).
%
% Figures (inline TikZ; promote to figures/ later if reused):
%   - fig:engine-pipeline; index construction + retrieval flow.
%
% Constraints:
%   - Describe, do not pitch. Studied subject, not contribution.
%   - Stay within ~2.5 pp; spill detail into Appendix if needed.
%   - No vendor names for backend infrastructure; tree-sitter only.
% =========================================================================

\section{System}
\label{sec:system}

This section describes the studied subject: the SuperCoder coding-agent harness,
the context engine that the ON/OFF arms ablate, and the OpenCode comparator that
the cross-harness arm runs.

\subsection{SuperCoder harness}
\label{sec:system-harness}

SuperCoder \citep{supercoder2024} is a coding agent built around a single-LLM control loop. The
shipped binary supports three modes (Ask, Plan, Coding); the evaluation runs in Coding
mode, the only mode that may write files or execute shell commands, so all
mechanics described below refer to the Coding-mode loop. A provider gateway sits
in front of the LLM client and captures token-level cost per turn uniformly across
arms (\S\ref{sec:design-sandbox}).

Each turn assembles a prompt (system instructions, tool schema, message history)
and issues a single LLM call. If the model emits tool calls, the harness
dispatches them, awaits results, and appends them to the message history; if it
emits text with no tool calls, the loop terminates. The reasoning-and-acting
posture follows the ReAct \citep{react2023} pattern, and the tool-call interface
follows the function-calling line introduced by Toolformer \citep{toolformer2023}.
Multiple tool calls in a single response are executed in parallel. If the rolling token count exceeds a
threshold, the harness compacts the message history by summarizing older turns;
the compaction step is disclosed for reproducibility but is not load-bearing for
the ablation. The loop terminates when (a) the model emits no tool calls,
(b) a configured per-cell turn budget is exhausted, or (c) the 30-minute
per-cell wall-clock cap (\S\ref{sec:design-sandbox}) elapses.

The agent calls a fixed tool set: \texttt{read}, \texttt{write}, \texttt{edit},
\texttt{bash}, \texttt{git}, \texttt{grep}, and \texttt{glob}, plus task-management
tools (\texttt{todo\_write}, \texttt{apply\_patch}). All of these are identical
across SC-ON and SC-OFF. The two context-engine tools, \texttt{codebase\_search}
and \texttt{codebase\_graph}, are available only in SC-ON; SC-OFF removes those
two tools from the schema and changes nothing else
(\S\ref{sec:system-engine}).

\subsection{Context engine}
\label{sec:system-engine}

The context engine is a separate service that the agent calls through two tools.
It maintains a per-repository index that is built once on first contact and
updated incrementally on subsequent runs via Merkle-tree diffs over the working
copy, so a source edit invalidates and re-indexes only the affected chunks
rather than the whole repository. Each cell in this study starts from a fresh
sandbox, so every arm exercises the build path; the incremental-update path is
part of the engine but not load-bearing in the run. The index has three
components: a \textbf{vector index} of code-chunk embeddings for semantic
similarity, a \textbf{graph index} of definitions and call edges for structural
reachability, and a \textbf{lexical (BM25) index} of identifiers and tokens for
exact-match recall. Index construction begins with \texttt{tree-sitter} parsing
per source file; definitions, references, and call edges are extracted from the
resulting AST and chunked for embedding.

Figure~\ref{fig:engine-pipeline} sketches the indexing pipeline and the retrieval
path. The components are named at the level the public eval repo can support
(\S\ref{sec:integrity-ledger}); the backend service that hosts the three indices
is internal and not part of the released artifact.

\begin{figure}[!tbp]
  \centering
  \resizebox{\linewidth}{!}{%
  \begin{tikzpicture}[
      x=1.3cm, y=1.15cm,
      every node/.style={font=\small},
      box/.style={draw, rounded corners=3pt, align=center,
                  inner sep=6pt, minimum height=13mm,
                  text width=36mm, font=\small},
      idx/.style={box, fill=black!5},
      tool/.style={box, fill=black!12},
      arr/.style={-Latex, thick, line width=0.7pt},
      phaselbl/.style={font=\small\bfseries\itshape, fill=white, inner sep=3pt},
    ]
    %% Phase 1 title
    \node[phaselbl] at (0, 5.0) {Indexing (build once, then Merkle-diff updates)};

    %% Indexing chain (vertical, centered)
    \node[box] (repo)    at (0, 4.0) {\textbf{Repository}\\\scriptsize source files};
    \node[box] (parse)   at (0, 2.7) {\textbf{\texttt{tree-sitter} parse}\\\scriptsize AST per source file};
    \node[box] (extract) at (0, 1.4) {\textbf{Symbol + call-graph extraction}\\\scriptsize definitions, identifiers, call edges};
    \node[box] (chunk)   at (0, 0.1) {\textbf{Chunking + embedding}\\\scriptsize code chunks $\to$ vectors};

    \draw[arr] (repo)    -- (parse);
    \draw[arr] (parse)   -- (extract);
    \draw[arr] (extract) -- (chunk);

    %% Three indices (shared band)
    \node[idx] (bm25)   at (-3.4,-1.6) {\textbf{Lexical (BM25)}\\\scriptsize identifier + token matches};
    \node[idx] (graph)  at ( 0,  -1.6) {\textbf{Call-graph index}\\\scriptsize caller / callee edges};
    \node[idx] (vector) at ( 3.4,-1.6) {\textbf{Vector index}\\\scriptsize semantic similarity};

    \draw[arr] (chunk.south) -- (bm25.north);
    \draw[arr] (chunk.south) -- (graph.north);
    \draw[arr] (chunk.south) -- (vector.north);

    %% Phase separator + Phase 2 title
    \draw[dashed, gray!70] (-5.0, -2.8) -- (5.0, -2.8);
    \node[phaselbl] at (0, -3.1) {Retrieval (per agent call)};

    %% Retrieval row (horizontal: tool -> hybrid -> results)
    \node[tool] (tools)   at (-3.4, -4.6) {\textbf{\texttt{codebase\_search} /}\\\textbf{\texttt{codebase\_graph}}\\\scriptsize agent issues query};
    \node[box]  (hybrid)  at ( 0,   -4.6) {\textbf{Hybrid retrieval}\\\scriptsize fuse, rerank, dedup};
    \node[box]  (results) at ( 3.4, -4.6) {\textbf{Ranked results}\\\scriptsize paths + snippets + scores};

    \draw[arr] (tools)  -- (hybrid);
    \draw[arr] (hybrid) -- (results);

    %% Indices feed into hybrid retrieval (downward fan-in)
    \draw[arr] (bm25.south)   -- (hybrid.north west);
    \draw[arr] (graph.south)  -- (hybrid.north);
    \draw[arr] (vector.south) -- (hybrid.north east);
  \end{tikzpicture}%
  }% close \resizebox
  \caption{Context-engine pipeline. The upper block runs once per repository on
    first contact and re-runs incrementally on subsequent contacts via
    Merkle-tree diffs over the working copy: \texttt{tree-sitter} produces an
    AST per source file, an extractor walks the ASTs to collect definitions,
    identifiers, and call edges, and code chunks are embedded into vectors;
    the result is three indices populated in parallel. The lower block runs on
    every agent call:
    \texttt{codebase\_search} or \texttt{codebase\_graph} dispatches a query to
    hybrid retrieval, which fuses hits across the three indices and returns a
    ranked result list to the agent. Per-tool input and result schemas are
    described in \S\ref{sec:system-engine}.}
  \label{fig:engine-pipeline}
\end{figure}

\textbf{Agent-facing tools.} \texttt{codebase\_search} takes a natural-language query plus a retrieval strategy (vector, lexical, graph, or hybrid) and returns a ranked list of code chunks, each carrying file path, snippet, relevance score, and the index that produced it; a local overlay drops paths the agent has deleted and flags stale ones. \texttt{codebase\_graph} takes a symbol and traverses the call-graph index, returning callers and callees grouped by direction, each carrying the defining file path and the distance from the query node.

\textbf{ON versus OFF concretely.} SC-ON exposes both
\texttt{codebase\_search} and \texttt{codebase\_graph} to the agent alongside
the file I/O, shell, and search tools listed in \S\ref{sec:system-harness};
SC-OFF removes exactly those two tools from the schema and leaves everything
else identical (the rest of the tool set, the model, the prompt template, the
sandbox, the scorer; \S\ref{sec:design-arms}). The toggle is the headless
runner's command-line flag for the engine endpoint: present invokes SC-ON,
absent invokes SC-OFF. The ablation surface is therefore ``with versus without
the two engine tools,'' and the resolve, localization, and cost consequences of
that surface live in \S\ref{sec:results-ablation}.

\subsection{Comparator: OpenCode}
\label{sec:system-opencode}

OpenCode \citep{opencode2025} is an open-source coding agent: a single-LLM control loop with parallel tool dispatch and a fixed tool set built around \texttt{rg} (ripgrep), \texttt{read}, \texttt{glob}, and \texttt{bash}; no structural codebase index, no embedding-based search, no precomputed call-graph. In our evaluation, OpenCode runs in the same per-task container as the two SuperCoder arms, with Claude Opus~4.7 and the same 30-minute wall-clock cap (\S\ref{sec:design-sandbox}). The headline cross-harness comparison is \S\ref{sec:results-crossharness}.
\FloatBarrier
